% Edgar Allan Poe, "The Tell-Tale Heart" (1843, public domain), first part only, typeset like an anthology page:
% line numbers in the margin every 5 lines, short introduction, glossary, questions. English, 2 pages.
% The introduction, glossary and questions are our own text.
\documentclass[11pt,a4paper]{article}
\usepackage[T1]{fontenc}
\usepackage[utf8]{inputenc}
\usepackage{ebgaramond}
\usepackage[scaled=0.9]{helvet}
\usepackage[english]{babel}
\usepackage[inner=2.9cm,outer=2.2cm,top=2.2cm,bottom=2.2cm]{geometry}
\usepackage{xcolor}
\usepackage{tcolorbox}
\usepackage{enumitem}
\usepackage{fancyhdr}
\usepackage{microtype}
\usepackage{lineno}

\definecolor{anthology}{HTML}{6B1E2E}
\setlength{\parindent}{1.3em}
\setlength{\parskip}{0pt}
\linenumbersep=1.1em
\renewcommand{\linenumberfont}{\sffamily\scriptsize\color{anthology}}
\modulolinenumbers[5]

\pagestyle{fancy}
\fancyhf{}
\renewcommand{\headrulewidth}{0.4pt}
\fancyhead[L]{\sffamily\small\textcolor{anthology}{\textbf{AMERICAN GOTHIC}} \textperiodcentered{} Short Fiction}
\fancyhead[R]{\sffamily\small Edgar Allan Poe}
\fancyfoot[L]{\sffamily\footnotesize Text: public domain (Poe, 1843). Introduction, glossary and questions: our own text, for testing.}
\fancyfoot[R]{\sffamily\footnotesize \thepage}

\begin{document}

% ---- Title ---------------------------------------------------------------------------------
\begin{center}
{\sffamily\small\color{anthology}\textbf{UNIT 3 \quad THE UNRELIABLE NARRATOR}}\\[0.9em]
{\fontsize{26}{30}\selectfont\textsc{The Tell-Tale Heart}}\\[0.5em]
{\large\itshape Edgar Allan Poe}\\[0.1em]
{\sffamily\small 1809--1849 \textperiodcentered{} first published in 1843}
\end{center}

\vspace{0.4em}
\begin{tcolorbox}[colback=anthology!6, colframe=anthology, boxrule=0.5pt, left=8pt, right=8pt, top=5pt, bottom=5pt,
  arc=0pt, title={\sffamily\bfseries Before you read}, fonttitle=\sffamily\small, coltitle=white, colbacktitle=anthology]
\sffamily\small
Edgar Allan Poe was an American writer of poems, tales and criticism, best known for his stories of fear and
guilt. ``The Tell-Tale Heart'' is told by a nameless narrator who tries to convince us that he is sane, while he
describes a murder he planned for a week. As you read, notice \textbf{what the narrator says about himself} and
what his words reveal that he does not intend to show. The excerpt below is the first part of the story; the
numbers in the margin count every fifth line.
\end{tcolorbox}

\vspace{1.0em}
% ---- Text with line numbers --------------------------------------------------------------------
\begin{linenumbers}
\noindent\textsc{True!}---nervous---very, very dreadfully nervous I had been and am; but why \emph{will} you say
that I am mad? The disease had sharpened my senses---not destroyed---not dulled them. Above all was the sense of
hearing acute. I heard all things in the heaven and in the earth. I heard many things in hell. How, then, am I
mad? Hearken! and observe how healthily---how calmly I can tell you the whole story.

It is impossible to say how first the idea entered my brain; but once conceived, it haunted me day and night.
Object there was none. Passion there was none. I loved the old man. He had never wronged me. He had never given
me insult. For his gold I had no desire. I think it was his eye! yes, it was this! He had the eye of a
vulture---a pale blue eye, with a film over it. Whenever it fell upon me, my blood ran cold; and so by
degrees---very gradually---I made up my mind to take the life of the old man, and thus rid myself of the eye
forever.

Now this is the point. You fancy me mad. Madmen know nothing. But you should have seen me. You should have seen
how wisely I proceeded---with what caution---with what foresight---with what dissimulation I went to work! I was
never kinder to the old man than during the whole week before I killed him. And every night, about midnight, I
turned the latch of his door and opened it---oh so gently! And then, when I had made an opening sufficient for my
head, I put in a dark lantern, all closed, closed, that no light shone out, and then I thrust in my head. Oh, you
would have laughed to see how cunningly I thrust it in! I moved it slowly---very, very slowly, so that I might
not disturb the old man's sleep. It took me an hour to place my whole head within the opening so far that I could
see him as he lay upon his bed. Ha!---would a madman have been so wise as this? And then, when my head was well in
the room, I undid the lantern cautiously---oh, so cautiously---cautiously (for the hinges creaked)---I undid it
just so much that a single thin ray fell upon the vulture eye. And this I did for seven long nights---every night
just at midnight---but I found the eye always closed; and so it was impossible to do the work; for it was not the
old man who vexed me, but his Evil Eye. And every morning, when the day broke, I went boldly into the chamber, and
spoke courageously to him, calling him by name in a hearty tone, and inquiring how he has passed the night. So you
see he would have been a very profound old man, indeed, to suspect that every night, just at twelve, I looked in
upon him while he slept.

Upon the eighth night I was more than usually cautious in opening the door. A watch's minute hand moves more
quickly than did mine. Never before that night had I felt the extent of my own powers---of my sagacity. I could
scarcely contain my feelings of triumph. To think that there I was, opening the door, little by little, and he not
even to dream of my secret deeds or thoughts. I fairly chuckled at the idea; and perhaps he heard me; for he moved
on the bed suddenly, as if startled. Now you may think that I drew back---but no. His room was as black as pitch
with the thick darkness, (for the shutters were close fastened, through fear of robbers,) and so I knew that he
could not see the opening of the door, and I kept pushing it on steadily, steadily.

I had my head in, and was about to open the lantern, when my thumb slipped upon the tin fastening, and the old man
sprang up in the bed, crying out---``Who's there?''

I kept quite still and said nothing. For a whole hour I did not move a muscle, and in the meantime I did not hear
him lie down. He was still sitting up in the bed listening;---just as I have done, night after night, hearkening to
the death watches in the wall.

Presently I heard a slight groan, and I knew it was the groan of mortal terror. It was not a groan of pain or of
grief---oh, no!---it was the low stifled sound that arises from the bottom of the soul when overcharged with awe.
I knew the sound well. Many a night, just at midnight, when all the world slept, it has welled up from my own
bosom, deepening, with its dreadful echo, the terrors that distracted me. I say I knew it well. I knew what the old
man felt, and pitied him, although I chuckled at heart. I knew that he had been lying awake ever since the first
slight noise, when he had turned in the bed. His fears had been ever since growing upon him. He had been trying to
fancy them causeless, but could not. He had been saying to himself---``It is nothing but the wind in the
chimney---it is only a mouse crossing the floor,'' or ``It is merely a cricket which has made a single chirp.''
Yes, he had been trying to comfort himself with these suppositions: but he had found all in vain. All in vain;
because Death, in approaching him had stalked with his black shadow before him, and the shadow had now reached and
enveloped the victim. And it was the mournful influence of the unperceived shadow that caused him to feel---although
he neither saw nor heard---to feel the presence of my head within the room.
\end{linenumbers}

\vspace{0.6em}
\begin{center}\sffamily\small\color{anthology}[ The story continues. The excerpt ends here. ]\end{center}

% ---- Glossary and questions ------------------------------------------------------------------
\vspace{0.4em}
\begin{tcolorbox}[colback=white, colframe=anthology, boxrule=0.5pt, left=8pt, right=8pt, top=5pt, bottom=5pt,
  arc=0pt, title={\sffamily\bfseries Glossary}, fonttitle=\sffamily\small, coltitle=white, colbacktitle=anthology]
\sffamily\small
\textbf{hearken}: to listen closely \quad \textbf{dissimulation}: hiding one's true feelings or plans \quad
\textbf{sagacity}: sharp judgement, wisdom \quad \textbf{death watches}: small beetles whose ticking in old wood
was said to warn of a death \quad \textbf{mortal}: connected with death \quad \textbf{suppositions}: guesses, things
one supposes to be true
\end{tcolorbox}

\vspace{0.6em}
\begin{tcolorbox}[colback=anthology!6, colframe=anthology, boxrule=0.5pt, left=8pt, right=8pt, top=5pt, bottom=5pt,
  arc=0pt, title={\sffamily\bfseries Think about it}, fonttitle=\sffamily\small, coltitle=white, colbacktitle=anthology]
\sffamily\small
\begin{enumerate}[leftmargin=1.5em, itemsep=0.2em, label=\textbf{\arabic*.}]
  \item The narrator opens with ``why will you say that I am mad?'' Which details in lines 1 to 10 show that he is
        trying to persuade a listener?
  \item He gives no reason for the murder except the old man's eye. Why might the author have chosen so strange a
        motive?
  \item Find three places where the narrator repeats a word or a phrase. What effect does the repetition have?
\end{enumerate}
\end{tcolorbox}

\end{document}
